% SPDX-License-Identifier: GPL-3.0-or-later
% Copyright (C) 2026 Alessandro Greco
\chapter{Выбор маршрута, входные узлы и vanguards}\label{cap:13-percorso}
\aggiornatoal{23 сентября 2026 г.}{коммит torspec \texttt{8e1f245}; tor 0.4.9.13; публичный консенсус от 23.09.2026, 15:00 UTC}

\begin{obiettivi}
К концу этой главы вы сможете:
\begin{itemize}[nosep]
  \item описывать порядок и ограничения, которые tor использует при выборе реле для цепи;
  \item объяснять, как пропускная способность в консенсусе и \emph{bandwidth-weights} распределяют нагрузку между позициями цепи;
  \item объяснять, зачем нужны входные узлы (guards), и прослеживать их эволюцию от трёх узлов к одному, а затем к двум;
  \item описывать алгоритм входных узлов: выборку, отфильтрованные, подтверждённые и основные узлы;
  \item объяснять, что такое vanguards и зачем они нужны onion-сервисам;
  \item моделировать выбор маршрута по реальному консенсусу публичной сети (лабораторная работа~7).
\end{itemize}
\end{obiettivi}

Система каталогов (\cref{cap:12-directory}) предоставляет клиенту список реле. Остаётся вопрос, который сильнее почти любого другого влияет на защиту пользователя: \emph{какие именно} реле следует использовать? Если противник контролирует одновременно первое и последнее реле цепи, он может коррелировать входящий и исходящий трафик и определить, кто с кем общается (\cref{cap:26-attacchi}). Поэтому выбор маршрута --- во многом искусство сделать такое событие как можно более редким, не жертвуя производительностью.

% ---------------------------------------------------------------------------
\section{Порядок и ограничения}\label{sec:percorso-vincoli}

Для трёхпереходной цепи tor сначала выбирает \textbf{выходной узел (exit)}, затем входной (guard) и лишь после этого промежуточный (middle)~\spec{path-spec}{Path selection and constraints}. Выходной узел выбирается первым, потому что именно к нему предъявляется больше всего ограничений: политика выхода должна разрешать запрошенный пункт назначения. Если клиент знает только имя назначения, но не его адрес, клиент выбирает выходной узел, который \emph{может} принять соединение: его политика не должна отклонять нужный порт до разрешающего правила.

За несколькими исключениями все маршруты удовлетворяют следующим ограничениям:
\begin{itemize}
  \item не использовать реле без флага \texttt{Fast};
  \item одно реле не может появляться дважды;
  \item не более одного реле из каждой IPv4-сети /16 (или IPv6 /32): два реле в одном дата-центре могут находиться под контролем одного оператора;
  \item не более одного реле из каждого \textbf{семейства}: операторы нескольких реле должны объявлять их как одно семейство, а tor избегает помещать два из них в одну цепь;
  \item первое реле является входным узлом (guard) (\cref{sec:guard});
  \item выходные узлы с флагом \texttt{BadExit} не используются.
\end{itemize}
Соединения, которые предположительно будут жить долго (tor относит к ним некоторые порты, например SSH 22 и IRC 6667), дополнительно требуют реле \texttt{Stable}. Onion-сервисы добавляют ограничение, не зависящее от порта: цепи введения на стороне сервиса и маршруты встречи строятся как цепи \texttt{Stable}, а их конечные точки не должны иметь \texttt{MiddleOnly}~\spec{path-spec}{Path selection and constraints}. Упомянутые \enquote{несколько исключений} в основном относятся к цепям onion-сервисов с vanguards (\cref{sec:vanguards}), где некоторые ограничения намеренно ослабляются, например входной узел может занимать промежуточную позицию. Запрет таких комбинаций сообщил бы противнику, какие реле \emph{не} были выбраны, и тем самым раскрыл бы сведения о выборе входных узлов.

\begin{approfondimento}[Два способа объявить семейство]
Традиционно оператор объявляет семейство, перечисляя все остальные свои реле в дескрипторе каждого реле. Два реле считаются принадлежащими одному семейству только если перечисляют друг друга. При большом числе реле списки становятся огромными, а одна ошибка конфигурации может нарушить связи. Новый механизм \emph{happy families} даёт каждому семейству ключ. Каждое реле доказывает принадлежность сертификатом, подписанным этим ключом (\texttt{family-cert} в дескрипторе, \texttt{family-ids} в микродескрипторе). Параметры консенсуса \texttt{use-family-lists} и \texttt{use-family-ids} определяют, какой механизм применяют клиенты.
\end{approfondimento}

% ---------------------------------------------------------------------------
\section{Взвешивание по пропускной способности}\label{sec:percorso-pesi}

\subsection{Зачем нужны веса}

Если бы клиенты выбирали реле равномерно случайно, реле на 1\,MB/s получало бы столько же трафика, сколько реле на 1\,GB/s. Первое было бы перегружено, второе недоиспользовано. Поэтому Tor выбирает каждое реле с вероятностью, \textbf{пропорциональной его пропускной способности} в консенсусе (строка \texttt{w}, \cref{sec:dir-anatomia}), распределяя нагрузку по доступной мощности.

\subsection{Позиции цепи неравноценны}

Есть осложнение. Реле с флагом \texttt{Exit} сравнительно мало, а их исходящая пропускная способность --- самый дефицитный ресурс сети. Использовать выходной узел в качестве входного или промежуточного означало бы расходовать ресурсы, необходимые для выхода из сети. В меньшей степени то же относится к входным узлам. Поэтому авторитетные серверы каталогов рассчитывают для каждого консенсуса \textbf{bandwidth-weights}: коэффициенты, умножающие пропускную способность реле в зависимости от его \emph{типа} (только входной, только выходной, оба типа или ни один из них) и \emph{позиции}, для которой его выбирают~\spec{path-spec}{Weighting node selection}.

Каждый вес называется \texttt{W}\emph{pt}: \emph{p} --- позиция (\texttt{g} для входного, \texttt{m} для промежуточного, \texttt{e} для выходного, \texttt{b} для запросов каталога), а \emph{t} --- тип (\texttt{g}, \texttt{e}, \texttt{d} для совмещения входной и выходной ролей, \texttt{m} при отсутствии обеих). Значения выражаются в десятитысячных долях: \texttt{Wmg=3805} означает, что при выборе входного узла в промежуточную позицию учитывается только 38,05\,\% его пропускной способности.

\subsection{Веса реальной сети}

Следующие веса взяты из публичного консенсуса от 23 сентября 2026 года, 15:00 UTC, загруженного из CollecTor в лабораторной работе~7:
\begin{center}
\begin{tabular}{lrrrr}
\toprule
Позиция $\backslash$ тип & guard (\texttt{g}) & middle (\texttt{m}) & exit (\texttt{e}) & guard+exit (\texttt{d}) \\
\midrule
guard (\texttt{Wg\,}) & 6195 & 6195 & --- & 0 \\
middle (\texttt{Wm\,}) & 3805 & 10000 & 0 & 0 \\
exit (\texttt{We\,}) & 10000 & 10000 & 10000 & 10000 \\
\bottomrule
\end{tabular}
\end{center}
Следствия очевидны. Выходные узлы и реле, имеющие одновременно роли входного и выходного узла, имеют \textbf{нулевой} вес во входной и промежуточной позициях. В этой сети пропускная способность выходных узлов (33\,\% общей) настолько ценна, что резервируется для выходной позиции. Входных узлов больше: на них приходится 70\,\% пропускной способности; их вес равен 62\,\% на входной позиции и 38\,\% на промежуточной. Как покажет лаборатория, в 60\,\% цепей промежуточный узел сам имеет флаг \texttt{Guard}.

Формулы, по которым управляющие каталогами вычисляют эти веса, приведены в спецификации каталогов. Они различают несколько случаев в зависимости от того, пропускная способность какого типа наиболее дефицитна, и по возможности стремятся выровнять доступную ёмкость трёх позиций.

% ---------------------------------------------------------------------------
\section{Входные узлы (guards)}\label{sec:guard}

\subsection{Какую проблему они решают}

Предположим, противник контролирует долю $c$ пропускной способности возможных первых и выходных реле. Если клиент выбирает новые входной и выходной узлы случайно для каждой цепи, вероятность, что оба принадлежат противнику, составляет приблизительно $c^2$ на цепь. После $C$ цепей вероятность хотя бы одной компрометации равна $1 - (1 - c^2)^C$ и стремится к 1~\spec{guard-spec}{Introduction and motivation}. Пользователь, запускающий Tor каждый день, строит тысячи цепей; рано или поздно противник почти наверняка увидит часть его трафика.

Идея \textbf{входных узлов (guards)}, называемых в исследовательской литературе \emph{helper nodes}, состоит в том, чтобы выбрать первое реле \emph{один раз} и долго его использовать. Если входной узел честен, ни одна цепь не скомпрометирована таким способом; если он принадлежит противнику, компрометированы многие. Вероятность компрометации \emph{хотя бы один раз} снижается с почти 1 примерно до $c$.

\begin{center}
\begin{tabular}{rcc}
\toprule
Число цепей & Без фиксированного узла & С фиксированным входным узлом \\
\midrule
1 & 0.0025 & 0.0025 \\
100 & 0.221 & 0.050 \\
\num{1000} & 0.918 & 0.050 \\
\num{10000} & 1.000 & 0.050 \\
\bottomrule
\end{tabular}
\end{center}
\noindent Таблица, рассчитанная в лабораторной работе~7 для $c = 5\,\%$, показывает компромисс. Входные узлы не исключают компрометацию; они превращают событие, которое со случайными первыми реле со временем стало бы почти неизбежным, в редкое событие, затрагивающее небольшую долю пользователей вместо почти всех. Модель намеренно упрощена и, например, не учитывает смену выбранного guard со временем.

\subsection{Три, один, два}

Входные узлы появились в 2006 году после того, как Øverlier и Syverson показали, как локализовать onion-сервис, многократно заставляя его подключаться к контролируемым противником реле~\cite{overlier2006locating}. Много лет клиенты использовали три входных узла, меняя их каждые два-три месяца. Последующие исследования показали, что три входных узла и частая ротация увеличивают экспозицию. \proposta{236} (2014) предложила один входной узел, сохраняемый девять-десять месяцев. Но один входной узел делает связность более хрупкой и облегчает принудительную замену. \proposta{291} (2018) перешла к \textbf{двум} входным узлам. Встроенное значение по умолчанию у tor по-прежнему равен одному (\file{DFLT\_N\_PRIMARY\_GUARDS\_TO\_USE} в \file{src/feature/client/entrynodes.h}), но публичный консенсус содержит \texttt{guard-n-primary-guards-to-use=2}, переопределяя его для всех клиентов. Это наглядный пример управления сетью через консенсус.

\subsection{Алгоритм}

Современный алгоритм, описанный в \proposta[Предложении]{271} и спецификации входных узлов, управляет несколькими наборами узлов~\spec{guard-spec}{The algorithm}:
\begin{description}
  \item[Выборка входных узлов] Набор случайно выбранных с учётом пропускной способности входных узлов, сохраняемый на диске. Он растёт только при необходимости до максимума 60 реле или 20\,\% числа входных узлов в консенсусе. Ограничение выборки мешает противнику заставить клиента последовательно попробовать каждый входной узел сети, пока тот не достигнет узла противника.
  \item[Отфильтрованные входные узлы] Узлы из выборки, всё ещё присутствующие в консенсусе и пригодные при текущей конфигурации, например при наличии сетевого экрана, разрешающего только некоторые порты.
  \item[Подтверждённые входные узлы] Узлы, реально использованные для пользовательского трафика, в порядке предпочтения.
  \item[Основные входные узлы] Обычно три узла из отфильтрованного списка: сначала подтверждённые и отфильтрованные, затем при необходимости ещё не подтверждённые, но допустимые. Клиент использует первый доступный основной входной узел или первые два согласно параметру консенсуса. Если они недоступны, клиент временно пробует другие узлы.
\end{description}
\begin{figure}[htbp]
\centering
% SPDX-License-Identifier: GPL-3.0-or-later
% Copyright (C) 2026 Alessandro Greco
%% Русское альтернативное описание EPUB.
\altfig{Схема выбора сторожевых узлов Tor, а не цепь вложенных множеств. Узлы с флагом Guard в текущем консенсусе могут быть добавлены в постоянную выборку; выборка иногда хранит узлы, отсутствующие в новом консенсусе. FILTERED и CONFIRMED — разные, потенциально пересекающиеся подмножества выборки. Основные узлы (обычно три) берут только из FILTERED: сначала подтверждённые, затем при необходимости остальные отфильтрованные. Значения по умолчанию: до 60 узлов и порог 20\%; удаление через 120 дней для неподтверждённых, 60 дней после подтверждения либо 20 дней отсутствия в консенсусе. Основной узел может временно быть недоступен.}
% Отношения множеств и списков сторожевых узлов
\begin{tikzpicture}[font=\sffamily\scriptsize,>=Stealth,
  guardbox/.style={draw=torviola,fill=torviola!6,rounded corners=4pt,align=center,minimum height=1.1cm},
  guardarrow/.style={-Stealth,thick,torgrigio}]
  \node[guardbox,draw=torgrigio,fill=torgrigio!6,text width=10.5cm] (cons) at (5.8,5.1) {\textbf{Узлы с флагом Guard}: кандидаты из консенсуса};
  \node[guardbox,text width=10.5cm] (sample) at (5.8,3.7) {\textbf{SAMPLED}: сохранены; до 60, порог 20\,\%};
  \node[guardbox,text width=4.4cm] (filtered) at (2.7,2.0) {\textbf{FILTERED}: в консенсусе, разрешены};
  \node[guardbox,text width=4.4cm] (confirmed) at (8.9,2.0) {\textbf{CONFIRMED}: ранее использованы};
  \node[guardbox,draw=torverde!60!black,fill=torverde!15,text width=10.5cm] (primary) at (5.8,0.2) {\textbf{PRIMARY}: обычно 3 из FILTERED};
  \draw[guardarrow] (cons) -- (sample);
  \draw[guardarrow] (sample.south west) to[out=235,in=90] (filtered.north);
  \draw[guardarrow] (sample.south east) to[out=305,in=90] (confirmed.north);
  \draw[guardarrow] (filtered.south) to[out=270,in=160] (primary.west);
  \draw[guardarrow] (confirmed.south) to[out=270,in=20] (primary.east);
  \node[align=center,text=torgrigio,font=\sffamily\tiny] at (5.8,-1.05) {сначала подтверждённые; сроки удаления 120/60/20 дней};
\end{tikzpicture}

\caption{Связи между списками сторожевых узлов Tor: постоянная выборка формируется из узлов консенсуса, а FILTERED и CONFIRMED — разные, потенциально пересекающиеся подмножества выборки. PRIMARY выбирают только из FILTERED, отдавая приоритет подтверждённым, а затем остальным допустимым узлам. Ограничение выборки снижает риск контакта с новыми узлами.}
\label{fig:guard}
\end{figure}

Срок пребывания входного узла в выборке зависит от его состояния. Входной узел, который \textbf{никогда не подтверждался} (\texttt{CONFIRMED\_ON\_DATE} не задан), удаляется через \texttt{GUARD\_LIFETIME}, то есть 120 дней после выборки. \textbf{Подтверждённый} входной узел, реально используемый клиентом, удаляется через \texttt{GUARD\_CONFIRMED\_MIN\_LIFETIME} = 60 дней после подтверждения. Он также удаляется, если отсутствует в консенсусе 20 дней. Сохраняемые временные метки (\texttt{ADDED\_ON\_DATE} и \texttt{CONFIRMED\_ON\_DATE}) случайно сдвигаются назад максимум на \texttt{GUARD\_LIFETIME}/10 = 12 дней. Поэтому файл состояния не раскрывает точный момент начала использования Tor, а фактический срок службы подтверждённого входного узла становится 48--60 дней. Таким образом, примерно раз в два месяца сменяется именно \emph{подтверждённый} входной узел, а не раз в четыре. Эта деталь важна для симуляции в \cref{cap:29-ricerca}.

\begin{attenzione}[Новый входной узел --- это сигнал]
Tor тщательно различает \emph{недоступный} входной узел и \emph{неудачную} цепь. Если бы клиент менял входные узлы при каждом отказе цепи, противник мог бы намеренно ломать цепи через честные входные узлы, пока клиент не выберет враждебный. Именно эту проблему исправило обновление безопасности tor 0.4.9.13 от 23 сентября 2026 года: теперь клиент относит отказ к входному узлу только когда не может установить соединение, а не при сбоях цепей, потоков или запросов каталога.
\end{attenzione}

% ---------------------------------------------------------------------------
\section{Vanguards}\label{sec:vanguards}

Входные узлы защищают \emph{первый} переход. Onion-сервисы (\cref{cap:14-onion-services}) сталкиваются с дополнительной проблемой: вредоносный клиент может заставлять сервис строить сколько угодно цепей к выбранным атакующим точкам встречи. Если промежуточный узел каждой цепи выбирается заново случайно, рано или поздно он окажется у противника, который увидит входной узел сервиса. Следующим шагом может стать атака на этот входной узел или получение его данных. \textbf{Vanguards} расширяют идею входных узлов на последующие переходы, также фиксируя реле второго уровня, а в некоторых конфигурациях и третьего.

\begin{description}
  \item[Vanguards-lite] (\proposta{333}, tor 0.4.7.1-alpha) Включён по умолчанию у клиентов и onion-сервисов. Использует четыре реле второго слоя, не сохраняемые на диск, каждое со случайным сроком службы. Здесь код и спецификация расходятся. tor 0.4.9.13 использует диапазон 1--12 дней, в среднем около недели, заданный \file{MIN\_SECOND\_GUARD\_LIFETIME} и \file{MAX\_SECOND\_GUARD\_LIFETIME} в \file{src/feature/client/entrynodes.c}. Параметры консенсуса \texttt{guard-hs-l2-lifetime-min} и \texttt{-max} могут переопределить значения, но ни одного из них не было в консенсусе 23 сентября 2026 года. Спецификация vanguards, напротив, всё ещё указывает 1--22 дня, в среднем две недели. Режим рассчитан на сервисы и пользователей, активных не дольше нескольких недель.
  \item[Full vanguards] Реле фиксируются на двух уровнях, втором и третьем переходах, с более длительными сроками; предназначено для долго работающих onion-сервисов высокого риска. В tor реализуется через дополнение \emph{vanguards}, задающее \opt{HSLayer2Nodes} и \opt{HSLayer3Nodes}. Arti реализует режим непосредственно, имея \texttt{Full} наряду с \texttt{Lite}.
\end{description}
В лабораторной сети (\cref{sec:rp-stem}) мы видели цепи с назначением \file{HS\_VANGUARDS}: заранее построенные цепи с vanguards, готовые к расширению до точки введения или встречи.

% ---------------------------------------------------------------------------
\section{Лабораторная работа: исследование публичной сети}\label{sec:lab-percorso}

\begin{laboratorio}[7]{Выбор маршрута с учётом пропускной способности}
\textbf{Цель.} Написать \file{labs/ru/код/torlab/path\_selection.py} и использовать его для анализа реального консенсуса публичной сети.\\
\textbf{Предварительные условия.} Лабораторная работа~6 (\file{directory.py}). Публичный консенсус, загруженный из CollecTor, публичного архива Tor Project (\url{https://collector.torproject.org/}). Эти данные публикуются специально для исследований; их загрузка не затрагивает пользователей сети.\\
\textbf{Оценочное время.} 2 часа.
\end{laboratorio}

\paragraph{Шаг 1: вес реле.}
\begin{lstlisting}[language=Python]
def _relay_type(r):  # g = только Guard, e = только Exit, d = оба, m = ни одного
    g, e = "Guard" in r.flag, "Exit" in r.flag and "BadExit" not in r.flag
    return {(1, 1): "d", (1, 0): "g", (0, 1): "e", (0, 0): "m"}[(g, e)]

def weight(r, position, weights):
    key = f"W{position[0]}{_relay_type(r)}          # например "Wmg"
    return (r.bandwidth or 0) * weights.get(key, 10000) / 10000
\end{lstlisting}

\paragraph{Шаг 2: выбор маршрута.} Класс \texttt{Selector} один раз вычисляет веса для каждой позиции. Ограничения применяются \emph{отклонением}: выбирается реле, и если оно нарушает ограничение, например совпадает с уже выбранным или находится в том же /16, выбирается другое.
\begin{lstlisting}[language=Python]
def path(self, rng, guard=None):
    e = self.choose("exit", [], rng)
    g = guard or self.choose("guard", [e], rng)
    m = self.choose("middle", [e, g], rng)
    return g, m, e
\end{lstlisting}
Первая версия, которая для каждого выбора фильтровала более девяти тысяч реле и заново считала веса и сети, строила \num{2000} маршрутов за 310 секунд. Эта версия строит \num{20000} за 7 секунд. В статистическом моделировании производительность не является мелкой деталью.

\paragraph{Шаг 3: публичная сеть в числах.} Для консенсуса от 23 сентября 2026 года, 15:00 UTC:
\begin{lstlisting}[numbers=none]
реле: 9358
пропускная способность с Guard: 70.1%    пропускная способность с Exit: 33.1%
guard : подходят 5306; первые 10 = 1.7%; первые 100 = 12.6%
middle: подходят 8961; первые 10 = 2.8%; первые 100 = 13.7%
exit  : подходят 3126; первые 10 = 2.3%; первые 100 = 13.7%
20000 маршрутов: gge 30.9%  ggd 29.4%  gme 20.4%  gmd 19.3%
\end{lstlisting}
Число \emph{подходящих} реле меньше числа с соответствующим флагом: \cref{cap:06-architettura} сообщает 3177 реле \texttt{Exit}, а здесь их 3126. Функция \texttt{eligible} также требует \texttt{Running}, \texttt{Valid} и \texttt{Fast}, исключая \texttt{BadExit} и \texttt{MiddleOnly}. Разница состоит из 48 выходных реле без \texttt{Fast} и 3 с \texttt{BadExit}.
Последняя строка --- выборка с фиксированным зерном генератора 7. Без фиксированного зерна генератора значения при каждом запуске меняются на несколько десятых процентного пункта, и два эксперимента нельзя корректно сравнивать. Буквенные коды обозначают типы guard, middle и exit в маршруте. Например, \texttt{gme} означает guard, middle без обоих флагов и выходной узел без флага Guard. Как и предсказывают веса, ни один маршрут не использует выходной узел на первой или второй позиции; в 60\,\% случаев промежуточный узел имеет флаг \texttt{Guard}. Концентрация невелика: десять наиболее часто выбираемых выходных узлов вместе дают 2,3\,\% выборов, а первые сто --- 13,7\,\%.

\paragraph{Шаг 4: входные узлы и противники.} Класс \texttt{Adversary} вычисляет таблицу из \cref{sec:guard}. Измените долю противника: остаётся ли эффект входных узлов очевидным при $c = 1\,\%$ или $c = 20\,\%$?

% ---------------------------------------------------------------------------
\section*{Итоги}
\begin{itemize}
  \item tor сначала выбирает выходной, затем входной и промежуточный узлы, избегая повторения одного реле и реле из одной подсети /16 или одного семейства.
  \item Каждое реле выбирается с вероятностью, пропорциональной пропускной способности, умноженной на вес, зависящий от типа реле и позиции цепи; в рассмотренной сети выходные узлы используются только на выходной позиции.
  \item Входные узлы превращают компрометацию, которая со случайными первыми реле со временем становится почти неизбежной, в редкое событие; сейчас клиенты используют два входных узла.
  \item Алгоритм выбора входных узлов поддерживает ограниченную выборку, подтверждённые и основные узлы и осторожно связывает отказы с конкретным входным узлом.
  \item Vanguards фиксируют также второй слой onion-сервисов; vanguards-lite включён по умолчанию.
  \item Параметры консенсуса могут переопределять значения программы по умолчанию.
\end{itemize}

% ---------------------------------------------------------------------------
\section*{Упражнения}

\begin{esercizio}[1]\label{es:13-ordine}
Почему tor выбирает выходной узел первым? Что могло бы пойти не так, если сначала выбрать входной узел?
\end{esercizio}

\begin{esercizio}[1]\label{es:13-peso}
Реле имеет флаги \texttt{Guard} и \texttt{Exit} и пропускная способность в консенсусе 20000. Каков его вес в каждой из трёх позиций цепи при весах реальной сети, приведённых в этой главе?
\end{esercizio}

\begin{esercizio}[2]\label{es:13-guard}
Повторите расчёт таблицы для входных узлов для $c = 1\,\%$ и $c = 20\,\%$. После какого числа цепей вероятность компрометации без фиксированных входных узлов превышает 50\,\% в каждом случае?
\end{esercizio}

\begin{esercizio}[2]\label{es:13-sybil}
Противник хочет повысить вероятность выбора его реле в качестве выходных узлов. Какие ресурсы ему потребуется вложить и какие механизмы --- измерение пропускной способности, семейства, ограничения на один IP-адрес и веса --- усложняют задачу?
\end{esercizio}

\begin{esercizio}[2]\label{es:13-campione}
Почему выборка входных узлов ограничена 60 реле или 20\,\% их общего числа? Опишите атаку, которую предотвращает это ограничение.
\end{esercizio}

\begin{esercizio}[3]\label{es:13-famiglie}
Расширьте \texttt{Selector}, добавив ограничение по семействам по строкам \texttt{family} и \texttt{family-ids} микродескрипторов. Загрузите из CollecTor микродескрипторы, соответствующие консенсусу, и измерьте, сколько выборов отклоняется из-за принадлежности к семействам.
\end{esercizio}

\begin{esercizio}[3]\label{es:13-operatori}
Сгруппируйте реле публичного консенсуса по заявленным семействам и вычислите долю вероятности выбора выходного узла для десяти крупнейших семейств. Сравните с долей десяти индивидуально наиболее часто выбираемых выходных узлов. Что можно заключить о рисках операторов, управляющих многими реле?
\end{esercizio}

% ---------------------------------------------------------------------------
\section*{Дополнительная литература}
\begin{itemize}
  \item \titolo{Tor Path Specification} и \titolo{Tor Guard Specification}~\cite{torspec}.
  \item L.~Øverlier и P.~Syverson, \titolo{Locating Hidden Servers}~\cite{overlier2006locating}: атака, послужившая причиной введения входных узлов.
  \item Предложения 236, 271, 291 и 333: переход от одного входного узла к двум, современный алгоритм и vanguards-lite.
  \item Спецификация vanguards (vanguards-spec), включая таблицы временных шкал Sybil-атак.
\end{itemize}
